\section{Amplifying a mixed correlation}\label{sec:amplification}

Starting from the fixed nonzero correlation profile in
Section~\ref{sec:labels}, we construct a nonnegative divisor weight
$D_B$ that preserves it. We then expand the weighted correlation so
that Cauchy--Schwarz removes $g_x$ and leaves a positive quadratic
energy in $F_x$. The divisor weight depends only on primes tending to
infinity with $B$; its positive mean and bounded second moment are the
properties that will preserve the profile.

Retain $W$, $\phi$, and
\[
 \beta_*=
 \left|\int_{(0,\infty)\times\whZ}\phi(t)W(t,w)\,dt\,dw\right|>0
\]
from Section~\ref{sec:labels}, and all parameters and weights from
Section~\ref{sec:arithmetic}. In particular, $\phi$ is real and
nonnegative. Fix real, nonnegative, nonzero functions
$\rho_0,\psi\in C_c^\infty((1,2))$. The regularity grid and its tolerance
$\tau$ are fixed throughout. A constant $C_*$ is also fixed whenever
$B$ tends to infinity. Every limit in $x$ below is taken first, along
the bad subsequence already selected in Section~\ref{sec:labels}.
Constants may depend on these fixed functions, the labels, and the
grid and tolerance; dependence on $C_*$ will be indicated when it
matters.

\subsection{The divisor weight and the correlation it preserves}

For an integer $n\ge1$, consider the two factorizations
\[
 n=am,\qquad n+1=cl.
\]
We select $c$ with $\log c/B$ in the support of $\rho_0$ and $a$ with
$a/(Tc)$ in the support of $\psi$. The weights $A(a),A(c)$ apply to
the chosen divisors and $K(m),K(l)$ to the two quotients. Define their
normalized divisor sum on the profinite integers by
\begin{equation}\label{eq:11}
 D_B(w)=\frac1B
 \sum_{\substack{a\mid w\\ c\mid w+1}}
 A(a)K(w/a)A(c)K((w+1)/c)
 \rho_0\left(\frac{\log c}{B}\right)
 \psi\left(\frac{a}{Tc}\right),
 \qquad w\in\whZ.
\end{equation}
For integer $w=n$, this is exactly the sum over the two factorizations
above. In $\whZ$, divisibility means membership in the image of
multiplication by the divisor. Multiplication by a positive integer is
injective on $\whZ$, so each quotient in the sum is well defined.

For fixed $B$, the coefficient support is finite:
$e^B<c<e^{2B}$ and $Tc<a<2Tc$. Every coefficient is less than
$e^{4B}$ for sufficiently large $B$, so every prime of a nonzero
coefficient weight belongs to $\mathcal P$. The modulus
$\prod_{p\in\mathcal P}p^2$ suffices to determine $D_B$: a coefficient
is squarefree, and testing whether $p$ divides its quotient requires
at most the residue modulo $p^2$. Thus $D_B$ is a nonnegative function
of finitely many high-prime residues.

Define $I_{1,x}$ by the weighted correlation
\begin{equation}\label{eq:amp-weighted-correlation}
 \frac{I_{1,x}}B
 =\frac1x\sum_{n\ge1}
   \phi(n/x)\overline{g_x(n)}F_x(n+1)D_B(n).
\end{equation}
For each fixed $B$, the function $\phi(t)D_B(w)$ is a compactly
supported continuous profile test. The convergence in
\eqref{eq:labels-profile-convergence} therefore gives
\begin{equation}\label{eq:amp-profile-limit}
 \lim_{x\to\infty}\frac{I_{1,x}}B
 =\int_{(0,\infty)\times\whZ}
   \phi(t)W(t,w)D_B(w)\,dt\,dw.
\end{equation}

The factor $1/B$ in \eqref{eq:11} is explained by the fair-split
identity \eqref{eq:arith-fair-split}. Each finite prime set contributes
a factor $B$ when its divisors are summed as a fair split. For two
independent products with law $\nu_{1/2}$, the first-moment calculation
below will show that the smooth size-and-ratio expectation is a positive
constant divided by $B$, up to a smaller error. The two factors $B$ and
the external division by $B$ then give a bounded positive mean in that
model. The next reduction compares this model with the actual divisor
weight and bounds the exceptional cases.

\subsection{The moment target and the independent split model}

\begin{lemma}[Moments of the divisor weight]\label{lem:amp-divisor-moments}
There exist $C_0>0$ and $d_0>0$ such that, for each fixed
$C_*\ge C_0$ and all sufficiently large $B$,
\begin{equation}\label{eq:12}
 \|D_B\|_{L^2(\whZ)}\ll_{C_*}1,
 \qquad
 \int_{\whZ}D_B(w)\,dw\ge d_0.
\end{equation}
The constant $d_0$ is independent of $C_*$ and $B$.  In addition,
$\int D_B\ll1$ with a constant independent of $C_*$.
\end{lemma}

We prove this lemma after reducing the actual divisor weight to an
independent model and establishing the concentration estimate needed for
its second moment. A single split supplies the positive mean. Two splits
of the same prime sets govern the second moment.

Let $S_1,S_2$ be independent random subsets of $\mathcal P$, each
formed by including every prime $p$ independently with probability
$1/p$. Conditional on these sets, split each $S_\nu$ by independent fair
coins into a coefficient subset $\mathcal C_\nu$ and a remaining subset
$\mathcal R_\nu$. Write
\[
 a=\prod_{p\in\mathcal C_1}p,\qquad
 c=\prod_{p\in\mathcal C_2}p,
\]
with empty products equal to one, and define
\begin{equation}\label{eq:13}
 \widetilde D_B(S_1,S_2)
 =B\,\E_{\mathrm{split}}\left[
  \rho_0\left(\frac{\log c}{B}\right)
  \psi\left(\frac{a}{Tc}\right)
  \ind_{\text{all four subsets are regular}}
  \,\middle|\,S_1,S_2\right].
\end{equation}
In particular, $0\le\widetilde D_B\ll B$ pointwise.

\begin{lemma}[Reduction to independent fair splits]
\label{lem:amp-site-model}
For each fixed $C_*\ge0$, as $B\to\infty$,
\[
 \int_{\whZ}D_B(w)^r\,dw=\E\widetilde D_B^r+o(1)
 \qquad(r=1,2).
\]
\end{lemma}

\begin{proof}
For Haar distributed $w$, let
\[
 S_1(w)=\{p\in\mathcal P:p\mid w\},
 \qquad
 S_2(w)=\{p\in\mathcal P:p\mid w+1\}.
\]
At each prime these are disjoint hits, each of probability $1/p$.
Their joint law differs by $O(1/p^2)$ in total variation from two
independent Bernoulli$(1/p)$ indicators. Independence across primes
therefore couples them to the independent sets $S_1,S_2$ above, with
total failure probability
\[
 O\left(\sum_{p>P_0}p^{-2}\right)=O(1/P_0).
\]
The probability that any $p\in\mathcal P$ has $p^2\mid w$ or
$p^2\mid w+1$ is also $O(1/P_0)$.

On the complement of these exceptional events, coefficient divisors
are subset products of the corresponding site sets, and the prime set
of each quotient is exactly the complementary subset. By
\eqref{eq:arith-fair-split}, a divisor and its complement then have
untruncated weight $B2^{-|S_\nu|}$. The truncated weight inserts
exactly the regularity indicators for those two subsets. Summing the
two fair splits shows that the value of \eqref{eq:11} is precisely
\eqref{eq:13} on this successful coupling.

We can discard the exceptional events in both moments. Any nonzero
representation in \eqref{eq:11} forces each entire site to have at most
$(1+2\tau)\ell$ distinct primes, by the total regularity cutoffs for
the coefficient and quotient. This remains true in the presence of
squares: their two prime sets still cover the site's distinct primes,
although they need not be disjoint. Hence there are at most
$2^{2(1+2\tau)\ell}$ coefficient choices across the two sites.
Combining this with \eqref{eq:8} gives, for example, the uniform bound
$D_B\ll B^4$. Together with $\widetilde D_B\ll B$, the coupling and
site-square errors change its first two moments by at most
$O(B^8/P_0)=o(1)$. This proves the comparison.
\end{proof}

\subsection{The two-split question and an addition-product estimate}

The second moment of \eqref{eq:13} involves two conditionally
independent fair splits of the same two site sets. Denote their
coefficient products by $(a^{(h)},c^{(h)})$, $h=1,2$, and their
coefficient and remaining prime sets by
$\mathcal C_\nu^{(h)},\mathcal R_\nu^{(h)}$, $\nu=1,2$.
Choose fixed closed intervals
$I_{\rho_0}\subset(1,2)$ and $I_\psi\subset(1,2)$ containing the
supports of $\rho_0$ and $\psi$, respectively. Let $\mathcal H_h$
be the size-and-ratio event
\[
 \frac{\log c^{(h)}}B\in I_{\rho_0},
 \qquad
 \frac{a^{(h)}}{Tc^{(h)}}\in I_\psi,
\]
and let $\mathcal E_h$ be $\mathcal H_h$ together with regularity
of its four subsets. Since the smooth weights are bounded and
nonnegative,
\begin{equation}\label{eq:amp-two-split-reduction}
 \E\widetilde D_B^2
 \ll B^2\Prob(\mathcal E_1\cap\mathcal E_2).
\end{equation}
It therefore suffices to prove the global probability bound
$\Prob(\mathcal E_1\cap\mathcal E_2)\ll_{C_*}B^{-2}$.

To prove this target, we will group the changes between the two splits
by their largest logarithmic scale $Y$. A further unconditional coupling
will replace the first coefficient and remaining classes by independent
prime processes, with its error paid before conditioning. In that
comparison model, the additions from a remaining set at primes with
$\log p\le Y$ are independent selections with probabilities $1/(4p)$,
independent also of the data above $Y$. After fixing the retained
coefficient primes and the additions at the other site, the second ratio
window confines the logarithm of the recipient's addition product to an
interval of fixed length. A new prime with logarithm above $Y/2$ also
forces that product logarithm above $Y/2$. The next estimate
gives the resulting $O(1/Y)$ probability uniformly in the interval's
location. The second-moment proof will combine it with the retained
first-remainder tail restrictions and the high assignment coins.
The product itself need not be bounded by the largest allowed prime.

\begin{lemma}[Addition-product concentration]\label{lem:amp-small-ball}
Fix $H>0$.  Write $L=\log P_0$.  For sufficiently large $B$, let
$L\le Y\le8B$, put $U=\min(Y,4B)$, and form a random product $Q_Y$ by
including each prime $P_0<p\le e^U$ independently with probability
$1/(4p)$.  Uniformly over all intervals $J\subset\R$ of length at
most $H$,
\[
 \Prob\bigl(\log Q_Y\in J\cap[Y/2,\infty)\bigr)
 \ll_H\frac1Y.
\]
The implicit constant is independent of $B$, $Y$, and $J$.
\end{lemma}

\begin{proof}
If the intersection is empty there is nothing to prove.  Otherwise
enclose it in an interval $[v,v+H]$ with $v\ge Y/2$, and set
$X=e^v$.  The probability of a squarefree product $n$ of the permitted
primes is
\[
 q_Y\frac{w_Y(n)}n,
 \qquad
 q_Y=\prod_{P_0<p\le e^U}\left(1-\frac1{4p}\right),
 \qquad
 w_Y(n)=\prod_{p\mid n}\frac{1/4}{1-1/(4p)}.
\]
Set $w_Y(n)=0$ for other integers.  Since $Y/2\le U\le Y$ and
$U\ge L$, Mertens' estimate gives
\[
 q_Y\ll (L/U)^{1/4}\ll (L/Y)^{1/4},
 \qquad q_Y\le1.
\]

Let $c_{\mathrm{s}}>0$ be an admissible exponent for the
one-dimensional upper sieve of Section~\ref{sec:arithmetic}, with a
remainder $O(N^{.8})$ on intervals of length $N$.  Choose a fixed
$c>0$ sufficiently small that $c\le1/4$ and $2c<c_{\mathrm{s}}$, and
put $Z=e^{cY}$.  Then $Z\le e^U$ and, since $X\ge e^{Y/2}$,
$Z\le X^{c_{\mathrm{s}}}$.  Dropping squarefreeness, the restriction on
prime factors greater than $Z$, and the reducing weights above $Z$
majorizes $w_Y(n)$ by
\[
 w(n)=
 \prod_{p\le\min(P_0,Z)}\ind_{p\nmid n}
 \prod_{P_0<p\le Z} a_p^{\ind_{p\mid n}},
 \qquad a_p=\frac1{4-1/p}<1.
\]
Apply the random-weight form of that sieve on an interval containing
$[X,e^H X]$ and of length comparable to $X$, with constants depending
only on $H$.  It yields
\[
 \sum_{X\le n\le e^H X}w_Y(n)
 \ll_H X\mathcal D+X^{.8},
\]
where
\[
 \mathcal D=
 \prod_{p\le\min(P_0,Z)}(1-1/p)
 \prod_{P_0<p\le Z}\left(1-\frac{1-a_p}{p}\right).
\]
If $Z\ge P_0$, the relation $1-a_p=3/4+O(1/p)$ and Mertens'
estimate imply
\[
 \mathcal D\ll_c
 \frac1L\left(\frac{L}{Y}\right)^{3/4}.
\]
Consequently $q_Y\mathcal D\ll_c1/Y$.  If $Z<P_0$, ordinary rough
exclusion instead gives $\mathcal D\ll1/\log Z=1/(cY)$, and the same
conclusion follows from $q_Y\le1$.  Finally,
\[
 q_Y X^{-.2}\le e^{-Y/10}\ll\frac1Y.
\]
Using $1/n\le1/X$ on the summation interval proves
\[
 \Prob(X\le Q_Y\le e^H X)
 \le\frac{q_Y}{X}\sum_{X\le n\le e^H X}w_Y(n)
 \ll_H\frac1Y,
\]
as required.
\end{proof}

\subsection{The first and second moments}

\begin{proof}[Proof of Lemma~\ref{lem:amp-divisor-moments}]
By Lemma~\ref{lem:amp-site-model}, it suffices to prove the asserted
mean and second-moment bounds for $\widetilde D_B$.

\paragraph{The first moment.}
Without the regularity indicator in \eqref{eq:13}, the products
$a,c$ are independent and each has law $\nu_{1/2}$.  Let
$f_B=f_{1/2,B}$ and $f(s)=c_{1/2}s^{-1/2}$ on the compact positive
log ranges under consideration.  For fixed $s=\log c/B$ on the
support of $\rho_0$, \eqref{eq:2} and partial summation give, uniformly
there,
\[
 B\,\E_a\psi\left(\frac{a}{Tc}\right)
 =\int_0^\infty
    \psi(r)f_B\left(s+\frac{\log T+\log r}{B}\right)
    \frac{dr}{r}+o(1).
\]
To justify the use of \eqref{eq:2}, the relevant interval in
$\log a/B$ has length $O(1/B)$; its absolute distribution-function
error is $O(B^{-80})$, and partial summation against the smooth
window, whose total variation is bounded, leaves an error
$o(1/B)$ before the displayed multiplication by $B$.
Since $\log T/B\to0$ and $f_B$ converges uniformly with derivatives
on these compact intervals, another application of
\eqref{eq:2} to $c$ shows that the untruncated first moment tends to
\[
 d_*=
 \left(\int_0^\infty\rho_0(s)f(s)^2\,ds\right)
 \left(\int_0^\infty\psi(r)\,\frac{dr}{r}\right)>0.
\]
The same calculation, or its upper-bound version for fixed intervals
containing the supports, shows that the size-and-ratio event
$\mathcal H_1$ has probability $O(1/B)$ under the independent
coefficient-product law.

We next control the regularity losses in this normalization.  Cover
the coefficient support by $O(B)$ dyadic-size boxes
$c\asymp X$, $a\asymp TX$.  All their logarithmic scales lie in a
fixed compact subinterval of the ranges for the arithmetic
estimates.  From \eqref{eq:4}, the joint product probability at
$(a,c)$ is bounded by
\[
 C\frac{A_0(a)A_0(c)}{B^2ac}.
\]
Within one such box $ac\asymp TX^2$.  The independent-coefficient
version of \eqref{eq:9} bounds the untruncated weighted sum where
either coefficient is nonregular by
\[
 O\bigl(TX^2\{o(1)+e^{-c_3C_*}\}\bigr).
\]
After the harmonic denominators, summing the $O(B)$ boxes, and
multiplying by the $B$ in \eqref{eq:13}, the coefficient-regularity
loss is therefore $o(1)+O(e^{-c_3C_*})$.

Conditional on a selected coefficient subset at a site, the
remaining indicators at primes not selected for that coefficient
are independent, with probabilities
\[
 \frac{1/(2p)}{1-1/(2p)}=\frac1{2p-1}.
\]
The selected coefficient primes have reciprocal sum $O(B/P_0)$,
uniformly on the coefficient support, because the logarithm of
their product is $O(B)$.  Thus the independent-indicator regularity
estimate of Section~\ref{sec:arithmetic} applies uniformly after
these primes are omitted.  The conditional chance of a remaining
subset being nonregular is
$o(1)+O(e^{-c_3C_*})$.  Multiplying by the bounded smooth weights
and the coefficient-event probability $O(1/B)$ shows that the
remaining-subset loss in \eqref{eq:13} is again
$o(1)+O(e^{-c_3C_*})$ after its outside factor $B$.

It follows that
\[
 \E\widetilde D_B=d_*+O(e^{-c_3C_*})+o(1),
\]
where the error constant in the exponential term does not depend
on $C_*$.  This proves the asserted positive lower bound after
choosing $C_0$ large, for instance with a fixed $d_0<d_*/2$ and
then allowing the negligible coupling error.  The upper bound
$\int D_B\ll1$ follows by omitting all regularity indicators in
the independent model.  Both bounds can use constants independent
of $C_*$.  For the lower bound one may also note that increasing
$C_*$ only relaxes the tail cutoffs.

\paragraph{The second moment.}
We use the two-split notation introduced before
Lemma~\ref{lem:amp-small-ball}. By \eqref{eq:amp-two-split-reduction},
the required bound is the global probability estimate stated there.

Compare the two assignments of each site prime.  If there is a
change, its largest logarithm lies in one of the disjoint
intervals $(Y/2,Y]$, where
\[
 Y=2^i\log P_0,\qquad 1\le i\le i_{\max}.
\]
Here $4B\le2^{i_{\max}}\log P_0<8B$ for large $B$.
Write $\mathcal A_Y$ for the event that this is the largest changed
interval.  Every change in it is either a move from remaining to
coefficient or a move in the reverse direction.  The joint law
and $\mathcal E_1\cap\mathcal E_2$ are invariant under exchanging
the splits.  Hence
\[
 \Prob(\mathcal E_1\cap\mathcal E_2\cap\mathcal A_Y)
 \le 2\sum_{\nu=1}^2
 \Prob(\mathcal E_1\cap\mathcal E_2\cap\mathcal A_Y
       \cap\mathcal U_{Y,\nu}),
\]
where $\mathcal U_{Y,\nu}$ requires at least one prime at site
$\nu$ with logarithm in $(Y/2,Y]$ to move from the first remaining
set to the second coefficient set.  This orientation inequality
has been taken in the original symmetric two-split law.

For the upper bound on its right side, couple the four first-split
sets $\mathcal C_\nu^{(1)},\mathcal R_\nu^{(1)}$ to four mutually
independent prime processes, each with inclusion probabilities
$1/(2p)$.  At a prime the two classes at one site were mutually
exclusive, so the cost of this coupling is $O(1/p^2)$, and the
total cost is $O(1/P_0)$.  Assign independent fair second-split
coins to each occurrence in these processes.  On the successful,
collision-free coupling this constructs exactly the original
second split.  On the exceptional event there may be two
occurrences of a prime, for which we still assign independent
coins and define coefficient products with multiplicity.  This
provides a convenient coupled model for upper bounds.  Even if
the coupling error is counted separately in every one of the
$O(\log B)$ intervals, its contribution after multiplication by
$B^2$ is $O(B^2\log B/P_0)=o(1)$.  No uniform coupling assertion
conditioned on a rare coefficient value is needed.

Work now in that independent model. Its first coefficient products have
independent $\nu_{1/2}$ laws, so the first-moment calculation gives
$\Prob(\mathcal H_1)=O(1/B)$. Fix first coefficient sets satisfying
$\mathcal H_1$ and the first coefficient regularity conditions.
For each term on the oriented right side,
we will retain the needed first-remainder restrictions while bounding
its second ratio condition. The sum of these conditional bounds over
the changed scales, together with the no-change contribution, must be
$O_{C_*}(1/B)$ uniformly in the fixed coefficient sets. The orientation
factor and coupling error remain outside this conditional calculation.

For a prime set $S$ let
$N_S(>Y)=\#\{p\in S:\log p>Y\}$, and put
\[
 q(Y)=.4\log(B/Y)-C_*.
\]
The first coefficient tail cutoffs give
$N_{\mathcal C_\nu^{(1)}}(>Y)\ge q(Y)$.
Expose the first remaining sets above $Y$, and retain the two
necessary tail restrictions
$N_{\mathcal R_\nu^{(1)}}(>Y)\ge q(Y)$.
Conditional on all these first high-prime sets, the chance that
none of their second assignments changes is exactly
\[
 2^{-\sum_{\nu=1}^2
       \{N_{\mathcal C_\nu^{(1)}}(>Y)
        +N_{\mathcal R_\nu^{(1)}}(>Y)\}}.
\]
On the retained restrictions this is at most
\[
 2^{-4q(Y)}
 =2^{4C_*}(Y/B)^{\alpha_0},
 \qquad \alpha_0=1.6\log2>1.
\]
If $q(Y)<0$ this upper bound exceeds one, which remains a valid
bound.  Averaging over the first remaining high-prime sets does
not increase it.

All first remaining restrictions below $Y$ may now be omitted.
For each site, a prime below this threshold is newly added to
the second coefficient exactly when it lies in the independent
first remaining process and its second coin selects the
coefficient.  These new-addition indicators are independent and
have probabilities $1/(4p)$.  The new additions at the two sites
are independent of one another, of the first coefficient sets
and their second coins, and of all the exposed variables above
$Y$.

Fix a recipient site $\nu$.  Condition on the retained
first-coefficient primes at both sites and on the additions at
the other site.  The second ratio condition in $\mathcal H_2$
then restricts the logarithm of the addition product at site
$\nu$ to an interval of length at most
\[
 H_\psi=\log\left(\frac{\sup I_\psi}{\inf I_\psi}\right).
\]
For example, for recipient site 1 the second coefficient
products have the form $a^{(2)}=a_{\mathrm{ret}}b_{\mathrm{new}}$
and $c^{(2)}=c_{\mathrm{ret}}c_{\mathrm{new}}$, and solving
$a^{(2)}/(Tc^{(2)})\in I_\psi$ gives precisely such an interval
for $\log b_{\mathrm{new}}$.  The same calculation with the
inequality inverted applies at site 2.  On
$\mathcal U_{Y,\nu}$ there is a new prime with logarithm greater
than $Y/2$, so this addition product also has logarithm at least
$Y/2$.  There are no new additions above $Y$ on
$\mathcal A_Y$.

We can drop the second coefficient size restriction and all
second regularity restrictions for this upper bound.
Lemma~\ref{lem:amp-small-ball} then bounds the conditional chance
of the remaining addition-product conditions by $O(1/Y)$,
uniformly in every value on which we have conditioned.
To make the conditioning order explicit, write $\mathcal C$ for the
fixed first coefficient sets, $G_Y$ for the retained first remaining
tail restrictions, and $H_Y$ for the event of no change above $Y$.
Let $\mathcal F$ include $\mathcal C$ and reveal all second coins of
the first coefficient sets, all first remaining data and second coins
above $Y$, and the additions at the other site.  The recipient
addition product $Q_\nu$ below $Y$ is independent of this information,
and the ratio condition specifies an $\mathcal F$-measurable interval
$J(\mathcal F)$.  For the oriented event under consideration, the
tower property gives
\[
\begin{aligned}
 &\Prob(\text{oriented event}\mid\mathcal C)\\
 &\quad\le \E\!\left[\ind_{G_Y}\ind_{H_Y}
   \Prob\!\left(\log Q_\nu\in J(\mathcal F)\cap[Y/2,\infty)
        \mid\mathcal F\right)\,\middle|\,\mathcal C\right]\\
 &\quad\ll \frac1Y\E[\ind_{G_Y}\ind_{H_Y}\mid\mathcal C]
 \le \frac{2^{-4q(Y)}}Y.
\end{aligned}
\]
In the last step the high second coins are averaged after the
uniform small-ball bound is applied; no conditional bound on
$\Prob(H_Y\mid\mathcal F)$ is asserted.
Thus, conditional on any good first coefficient sets, the
contribution for this orientation at scale $Y$ is at most
\[
 O_{C_*}\left(\frac{(Y/B)^{\alpha_0}}{Y}\right).
\]
The constants here may depend on the fixed $C_*$, but not on
the first coefficient sets, $B$, or $Y$.

If there is no change at all, use the tail cutoff at
$Y=\log P_0$.  All four first subsets then have at least
$.4\log R-C_*$ primes.  The same calculation bounds the
conditional probability, retaining the first remaining tail
restrictions, by
\[
 O_{C_*}(R^{-\alpha_0}).
\]
The dyadic scales satisfy
\[
 \sum_Y\frac{(Y/B)^{\alpha_0}}Y
 =B^{-\alpha_0}\sum_Y Y^{\alpha_0-1}
 \ll \frac1B,
\]
because $\alpha_0>1$ and the largest $Y$ is less than $8B$.
Moreover,
\[
 B R^{-\alpha_0}
 = B^{1-\alpha_0}(1000\log B)^{\alpha_0}\longrightarrow0.
\]
Requiring first coefficient regularity only decreases the
$O(1/B)$ mass of $\mathcal H_1$ established above. Combining the preceding bounds,
including the orientation factor and the coupling errors,
gives
\[
 \Prob(\mathcal E_1\cap\mathcal E_2)
 \ll_{C_*}\frac1B\left(\frac1B+R^{-\alpha_0}\right)
      +O(\log B/P_0)
 \ll_{C_*}\frac1{B^2}+O(\log B/P_0).
\]
Multiplication by $B^2$ proves
$\E\widetilde D_B^2\ll_{C_*}1$. Lemma~\ref{lem:amp-site-model} transfers
these bounds to the Haar divisor weight and completes the proof of
\eqref{eq:12}.
\end{proof}

\subsection{The correlation survives the divisor weights}

\begin{lemma}[Preservation of the profinite profile]\label{lem:preservation}
For each fixed $C_*\ge C_0$, put $H(t,w)=\phi(t)W(t,w)$.  Then
\begin{equation}\label{eq:amp-preservation}
 \int H(t,w)D_B(w)\,dt\,dw
 -\left(\int H\,dt\,dw\right)\left(\int D_B\,dw\right)
 \longrightarrow0\qquad(B\to\infty).
\end{equation}
\end{lemma}

\begin{proof}
The function $H$ is bounded and supported on a fixed compact interval
in $t$, so it is square-integrable.
Conditional expectations with respect to $t$ and a finite
residue coordinate $w\bmod q$ approximate $H$ in $L^2$ as
these residue coordinates increase.  Equivalently, for every
$\varepsilon>0$ there is a function $H_q(t,w)$ depending only
on $t$ and $w\bmod q$ such that
\[
 \|H-H_q\|_2<\varepsilon,
 \qquad
 \int H_q\,dt\,dw=\int H\,dt\,dw.
\]
This follows, for example, by first approximating in the
dense space of finite sums of products of $L^2$ functions of
$t$ and locally constant functions of $w$, and then taking
conditional expectation.

For all sufficiently large $B$, no prime dividing $q$ belongs
to $\mathcal P$.  The finite-residue function $D_B$ is
therefore independent of $w\bmod q$ under Haar measure.
Consequently
\[
 \int H_q(t,w)D_B(w)\,dt\,dw
 =\left(\int H\,dt\,dw\right)\left(\int D_B\,dw\right).
\]
By Cauchy--Schwarz and \eqref{eq:12}, the error on replacing
$H_q$ by $H$ is $O_{C_*}(\varepsilon)$, uniformly in large
$B$; the length of the fixed $t$-support is absorbed in the
constant.  Since $\varepsilon$ is arbitrary, this proves
\eqref{eq:amp-preservation}.
\end{proof}

The factor $\int D_B\,dw$ in \eqref{eq:amp-preservation} is real,
nonnegative, and at
least $d_0$ eventually. Taking complex absolute values in
\eqref{eq:amp-profile-limit} and using \eqref{eq:amp-preservation}
therefore gives
\begin{equation}\label{eq:14}
 \liminf_{B\to\infty}\ \lim_{x\to\infty}
       \frac{|I_{1,x}|}{B}\ge d_0\beta_*.
\end{equation}
The possible dependence of the $L^2$ constant on $C_*$ is
harmless: this approximation is performed with $C_*$ fixed,
and the positive lower constant in \eqref{eq:14} is independent
of it.

\subsection{A positive quadratic energy and its diagonal}

The preserved correlation now has an expansion suited to
Cauchy--Schwarz. For $c\mid am+1$, put $l_a=(am+1)/c$. At fixed $B$,
the coefficient list in \eqref{eq:11} is finite. Thus
$g_x(am)=g_x(m)$ holds simultaneously for every coefficient $a$ on
that list once $x$ is sufficiently large, by
\eqref{eq:labels-multiplier}. Expanding
\eqref{eq:amp-weighted-correlation}, writing $n=am$, and using this
fixed-multiplier invariance gives the exact identity
\begin{equation}\label{eq:10}
 I_{1,x}=\frac1x\sum_{c\ge1}
   \rho_0\left(\frac{\log c}{B}\right)A(c)
   \sum_{m\ge1}K(m)\overline{g_x(m)}
   \sum_{\substack{a\ge1\\ c\mid am+1}}
   A(a)K(l_a)\psi\left(\frac{a}{Tc}\right)
   \phi\left(\frac{am}{x}\right)F_x(am+1).
\end{equation}
The values of $F_x$ are unchanged in this reindexing. The expanded
form places $g_x(m)$ outside the inner sum, where $|g_x(m)|=1$ supplies
the unit factor in Cauchy--Schwarz.

\begin{proposition}\label{prop:amp-energy}
There is $c_5>0$, independent of all sufficiently large fixed
$C_*$, for which the nonnegative energy below satisfies the displayed
lower bound:
\begin{equation}\label{eq:15}
\begin{gathered}
 I_{2,x}
 = \frac1x\sum_{c\ge1}
   \rho_0(\log c/B)A(c)
   \sum_{m\ge1}K(m)
   \left|
     \sum_{\substack{a\ge1\\c\mid am+1}}
     A(a)K(l_a)\psi(a/(Tc))
     \phi(am/x)F_x(am+1)
   \right|^2,
 \\
 \liminf_{B\to\infty}\ \liminf_{x\to\infty}
       \frac{I_{2,x}}{BT}\ge c_5
\end{gathered}
\end{equation}
Its diagonal contribution is negligible:
\[
 \lim_{B\to\infty}\ \limsup_{x\to\infty}
       \frac{I_{2,x}^{\mathrm{diag}}}{BT}=0.
\]
Here $I_{2,x}^{\mathrm{diag}}$ denotes the terms with equal
coefficient indices on expanding the square, and
$l_a=(am+1)/c$ as in \eqref{eq:10}.
\end{proposition}

\begin{proof}
Choose $0<\alpha_\phi<\beta_\phi<\infty$ with
$\operatorname{supp}\phi\subset[\alpha_\phi,\beta_\phi]$,
and write $u_- =\inf I_\psi$, $u_+=\sup I_\psi$.
If the inner sum in \eqref{eq:10} is nonzero, then
\[
 \frac{\alpha_\phi x}{u_+Tc}
 \le m\le
 \frac{\beta_\phi x}{u_-Tc}.
\]
Let $\mathcal I_{c,x}$ denote this interval. Use the nonnegative weights
$x^{-1}\rho_0(\log c/B)A(c)K(m)$ on this interval.
Cauchy--Schwarz gives
\[
 |I_{1,x}|^2
 \le V_{B,x} I_{2,x},
 \qquad
 V_{B,x}=
 \frac1x\sum_c\rho_0\left(\frac{\log c}{B}\right)A(c)
       \sum_{m\in\mathcal I_{c,x}}K(m),
\]
where $|g_x(m)|=1$ supplies the unit factor.  No multiplicativity of
the centered function $F_x$ is needed.

For fixed $B$, the weight $K_0$ is periodic and has mean
\[
 \int_{\whZ}K_0(w)\,dw
 =R^{1/2}\prod_{p\in\mathcal P}\left(1-\frac1{2p}\right)
 \ll1.
\]
Progression counting, $K\le K_0$, and the length of
$\mathcal I_{c,x}$ therefore imply
\[
 \limsup_{x\to\infty}V_{B,x}
 \ll\frac1T\sum_c
       \rho_0\left(\frac{\log c}{B}\right)\frac{A_0(c)}c
 \ll\frac BT.
\]
For the last bound, partition the compact log support of
$\rho_0$ into $O(B)$ dyadic-size boxes and apply
\eqref{eq:5} in each box.  Thus we may fix $C_4>0$,
independently of $C_*$, such that
\[
 \limsup_{x\to\infty}V_{B,x}\le C_4B/T
\]
for all sufficiently large $B$.  Combining this with
\eqref{eq:14} proves \eqref{eq:15}, with any fixed
\[
 0<c_5<\frac{(d_0\beta_*)^2}{C_4}.
\]
This makes the uniformity of the positive lower constant
explicit.

For the diagonal, set
\[
 Q_B=\sup A(a)K(l),
\]
where $a$ ranges over the coefficient support and $l$ over
positive integers for which the truncated weights are
nonzero.  The total-count cutoff and \eqref{eq:8} give
\[
 Q_B\ll B^{2(h_0+\tau\log2)+o(1)}=o(T),
\]
using the strict inequality
$2(h_0+\tau\log2)<.32$.  Since all smooth weights here are
nonnegative,
\[
 \psi(a/(Tc))^2\phi(am/x)^2
 \le \|\psi\|_\infty\|\phi\|_\infty
       \psi(a/(Tc))\phi(am/x).
\]
Bounding one of the two factors $A(a)K(l_a)$ by $Q_B$ and
$|F_x(am+1)|^2$ by 4 yields
\[
 I_{2,x}^{\mathrm{diag}}
 \ll Q_B\,\mathcal J_{B,x},
\]
where
\[
 \mathcal J_{B,x}=
 \frac1x\sum_c\rho_0\left(\frac{\log c}{B}\right)A(c)
 \sum_{m\ge1}K(m)
 \sum_{\substack{a\ge1\\c\mid am+1}}
 A(a)K(l_a)\psi\left(\frac{a}{Tc}\right)
 \phi\left(\frac{am}{x}\right).
\]
The same divisor reindexing as before gives
\[
 \mathcal J_{B,x}
 =\frac Bx\sum_{n\ge1}\phi(n/x)D_B(n).
\]
Here no label is present, so finite progression counting
directly gives
\[
 \lim_{x\to\infty}\mathcal J_{B,x}
 =B\left(\int_0^\infty\phi(t)\,dt\right)
       \left(\int_{\whZ}D_B(w)\,dw\right)
 \ll B
\]
by the first-moment upper bound in
Lemma~\ref{lem:amp-divisor-moments}.  It follows that
\[
 \limsup_{x\to\infty}\frac{I_{2,x}^{\mathrm{diag}}}{BT}
 \ll\frac{Q_B}{T}\longrightarrow0,
\]
which completes the proof.
\end{proof}
